\section{Introduction}\label{sec:intro}

Deep sequence models, gradient-boosted ensembles and, most recently, time-series foundation models are
increasingly used to forecast financial returns. Three problems recur in this literature. The first is
\emph{evaluation}: forecasts of price levels are scored by root mean squared error, which a random walk
nearly minimises, and directional accuracy is computed in ways that measure properties of the data
rather than of the forecaster. The second is \emph{opacity}: when a network does forecast well, it is
rarely possible to say what it has learned in terms a risk manager, a regulator or an economist would
recognise, let alone test. The third is \emph{leakage}: complex pipelines make it easy for information
from the forecast period to enter the forecast, and the resulting gains are both large and fictitious.
This paper is a response to all three, and it began as a response to the third: the authors'
own undergraduate thesis on a heterogeneous mixture-of-experts network for the CSI~300 contained two
look-ahead leaks and several illustrative figures that were not model output (\cref{app:leakage}).
Rebuilding that work from first principles produced the framework described here.

We make five contributions.

\paragraph{1. A formal account of the ``lagging paradox''.}
We show that the MSE-optimal forecast of tomorrow's price is today's price multiplied by the
conditional moment generating function of the return, so any forecaster trained on price levels
converges to a lagged random walk (\cref{prop:optimal-level}). The achievable reduction in price-level
MSE relative to the random walk is bounded by $s^2/(1+s^2)$, where $s$ is the daily
drift-to-noise ratio (\cref{prop:bound}); for any plausible $s$ this rules out RMSE gains of more than
about half a percent. A common directional-accuracy metric equals return sign persistence whenever the
forecaster lags (\cref{prop:da}). On nineteen assets a 60-lag price-level regression indeed collapses
to the random walk (coefficients summing to $0.994$--$1.003$) and loses to it in every market.

\paragraph{2. A white-box mixture of experts for conditional return distributions.}
The Semantic Mixture of Experts (SMoE) forecasts the full conditional distribution of the next-day
return as a mixture of linear-Gaussian experts with log-linear variances, weighted by a multinomial
logit gate in observable state variables. Researchers assign each expert a \emph{semantic group} of
covariates (trend, short-horizon, crisis) which acts as an exclusion restriction and identifies the
experts (\cref{prop:ident}). We derive conditions for geometric ergodicity, an ECM algorithm with
monotone likelihood, exact scores via Fisher's identity, QMLE asymptotics with sandwich inference, and
closed forms for every predictive functional: mean, variance decomposed into within-regime and
regime-disagreement risk, probability of a rise, VaR, ES and CRPS. Each forecast admits an exact additive
attribution to covariates and experts, and each gate variable an exact elasticity.

\paragraph{3. A yardstick for black boxes.}
Because the SMoE is a proper likelihood model, any black-box density forecaster can be projected onto it
by minimising expected Kullback--Leibler divergence. The resulting \emph{shadow} is a white-box
description of what the black box does, and the fidelity ratio $\Phi$ (the share of the black box's
out-of-sample log-score gain reproduced by its shadow) has an exact information-theoretic
interpretation (\cref{prop:fidelity}). We apply the yardstick to a deep heterogeneous MoE (the
undergraduate architecture, rebuilt as a mixture density network), to gradient boosting, to a
Markov-switching model and to three zero-shot foundation models (two sizes of Chronos-Bolt and
TimesFM~2.5). Because a pretrained model may have seen the evaluation period, we also propose a
release-date test for look-ahead through pretraining (\cref{prop:release}).

\paragraph{4. Specification and measurement tests for the white box.}
We derive a score test for omitted gate variables that needs only the restricted fit and reads as
a test of whether a candidate state variable predicts the regime surprise $\tau_{tk}-\pi_{tk}$
(\cref{prop:lm-gate}), and we show how measurement error in a daily variance estimator moves weight
from the daily to the weekly and monthly components of a log-HAR regression (\cref{prop:eiv}). Both
results are put to work: fourteen specifications of the experts and the gate are compared on all
nineteen series, and range-based variance covariates are compared with realised measures from
5-minute data, where the predicted attenuation of the daily loading matches the estimate.

\paragraph{5. Evidence from nineteen markets with a coverage guarantee and a deflated economic test.}
Using a leak-tested walk-forward protocol with a lockbox period, we evaluate thirteen forecasters
(including two sizes of Chronos-Bolt and TimesFM~2.5) on Chinese and international equity indices,
commodity and bond futures and Bitcoin. The SMoE has the best average log-score rank (2.95), is in
the 10\% model confidence set for 15 of 19 series on log score and 18 of 19 on volatility and tail
losses, and its 1\% VaR passes the Kupiec test on 17 of 19 series (GARCH-$t$: 10). It does not
forecast means or directions better than simple benchmarks: the hit rate averages about 51\%. Its white box
reveals that volatility, not momentum or drawdown, drives regime weights in the CSI~300, and that the
short-horizon expert exhibits continuation in Chinese indices but reversal in developed markets. An
adaptive conformal layer delivers distribution-free long-run VaR coverage, which combined with the SMoE
also removes violation clustering. Three experts forecast better out of sample than two, four or five,
and adding the VIX to the gate is the only change of specification that helps; with realised
variance from 5-minute data, a single-regime Realized GARCH forecasts the CSI~300 better in log score
than any SMoE variant, which we report as a limit of the regime structure. Trading the CSI~300 ETF from the SMoE predictive
distribution cuts the maximum drawdown from $-46\%$ to $-26\%$ (on index futures, from $-45\%$ to
$-38\%$), but after deflating for the strategies tried no excess return is statistically reliable.
In a lockbox period (July 2024 to October 2026) that we evaluated once, after every modelling choice
had been frozen, the mixture models again hold the top four log-score ranks among thirteen, with the
SMoE second behind LSTAR; GARCH-$t$ catches up on calibration (both pass the 1\% Kupiec test on 18 of
19 series);
there is still no reliable trading alpha; and the drawdown reduction is smaller (from $-19\%$ to
$-17\%$).

\paragraph{Organisation.}
\Cref{sec:literature} reviews related work. \Cref{sec:lagging} formalises the lagging paradox.
\Cref{sec:model,sec:estimation,sec:whitebox} develop the SMoE, its estimation and its white-box
interpretation, and \cref{sec:deployment} the calibration and decision layers.
\Cref{sec:montecarlo} reports simulations, \cref{sec:data} the data and protocol, and
\cref{sec:results-china,sec:results-global,sec:kgate,sec:rv,sec:yardstick,sec:economic,sec:lockbox}
the empirical results.

\Cref{sec:agenda} sets out a research agenda and \cref{sec:conclusion} concludes.
